\documentclass[../DoAn.tex]{subfiles}
\begin{document}

Chương 2 đã phân tích chi tiết yêu cầu chức năng, phi chức năng, mô hình hóa luồng nghiệp vụ bằng hệ thống biểu đồ UML và thiết kế cấu trúc cơ sở dữ liệu cho bài toán quản lý ra vào trạm. Chương 3 này sẽ cung cấp hệ thống nền tảng lý thuyết khoa học và mô tả chi tiết các giải pháp công nghệ được lựa chọn áp dụng trong đồ án tốt nghiệp. Việc xây dựng chương này là hoàn toàn cần thiết nhằm làm rõ cơ sở lý thuyết toán học của các thuật toán nhận diện và tính hợp pháp, an toàn của giao thức đọc thẻ định danh, bảo chứng cho tính khả thi của hệ thống trước khi cài đặt phần mềm thực tế. 

Các nội dung công nghệ cốt lõi trong chương được tổ chức thông qua bảy đề mục lớn bao gồm: (i) Lý thuyết học sâu nhận diện khuôn mặt metric learning với hàm mất mất ArcFace và kiến trúc mạng FaceNet tại mục 3.1, (ii) Tiêu chuẩn quốc tế ICAO Doc 9303 quy định cấu trúc dữ liệu chip NFC trên thẻ Căn cước công dân và cơ chế bắt tay bảo mật BAC tại mục 3.2, (iii) Giao thức mạng nhẹ MQTT tại mục 3.3, (iv) Kiến trúc hướng dịch vụ REST API tại mục 3.4, (v) Giải pháp xác thực JSON Web Token tại mục 3.5, (vi) Giao thức truyền thông thời gian thực WebSocket tại mục 3.6, và (vii) Hệ thống các framework lập trình chính được tích hợp sử dụng tại mục 3.7.

\section{Nhận diện khuôn mặt và không gian đặc trưng}
\label{section:3.1}

\subsection{Nền tảng lý thuyết trích xuất trắc lượng học}

Nhận diện khuôn mặt là bài toán nhận dạng mẫu cốt lõi trong lĩnh vực thị giác máy tính, hướng đến mục tiêu xác định danh tính của một cá nhân dựa trên hình ảnh kỹ thuật số của khuôn mặt. Trong hướng tiếp cận hiện đại dựa trên mạng nơ-ron tích chập sâu (Deep Convolutional Neural Network - DCNN), đặc trưng hình học khuôn mặt được biểu diễn dưới dạng vector nhúng có số chiều cố định gồm 512 chiều thực. 

Hai khuôn mặt thuộc cùng một người sẽ sinh ra hai vector nhúng có khoảng cách không gian nhỏ, ngược lại hai khuôn mặt khác người sẽ tạo ra khoảng cách lớn. Độ tương đồng Cosine (\textit{Cosine similarity}) giữa hai thực thể vector đặc trưng $\mathbf{A}$ và $\mathbf{B}$ được định nghĩa chính xác bởi Phương trình (\ref{eq:cosine}):

\begin{equation}
\label{eq:cosine}
    \text{similarity}(\mathbf{A}, \mathbf{B}) = \frac{\mathbf{A} \cdot \mathbf{B}}{\|\mathbf{A}\| \times \|\mathbf{B}\|}
\end{equation}

Phương trình (\ref{eq:cosine}) cho kết quả nằm trong không gian khoảng $[-1, 1]$, giá trị càng tịnh tiến về ngưỡng 1 thể hiện độ tương đồng hình học càng cao. Để phân định ranh giới giữa hai trạng thái, hệ thống thiết lập một giá trị ngưỡng quyết định bảo mật cố định làm căn cứ phân loại sinh trắc học.

\subsection{Biện luận giải pháp kỹ thuật lựa chọn}

Mô hình toán học này được lựa chọn nhằm giải quyết trực tiếp yêu cầu đối soát danh tính tại chỗ (Lớp xác thực 1 tại ứng dụng di động) và đối soát chéo an toàn (Lớp xác thực 2 tại máy chủ AI) đã đặt ra tại phân đoạn mục 2.1. Đề tài tiến hành đánh giá ưu nhược điểm giữa các hướng tiếp cận thuật toán khác nhau dựa trên hạ tầng thực tế triển khai nghiệp vụ của hệ thống.

\begin{table}[ht!]
\centering
\small
\renewcommand{\arraystretch}{1.3}
\caption{Bảng đối chiếu so sánh các thuật toán nhận diện khuôn mặt}
\label{tab:face_comparison}
\begin{tabular}{|>{\raggedright\arraybackslash}p{2.3cm}|>{\raggedright\arraybackslash}p{4.5cm}|>{\raggedright\arraybackslash}p{5.5cm}|>{\centering\arraybackslash}p{1.7cm}|}
\hline
\textbf{Thuật toán} & \textbf{Cơ chế cốt lõi} & \textbf{Ưu điểm / Hạn chế} & \textbf{Lựa chọn} \\ \hline
Eigenfaces & Phân tích thành phần chính PCA để trích xuất không gian đặc trưng tuyến tính. & Tính toán nhanh, tài nguyên thấp / Cực kỳ nhạy cảm với ánh sáng, góc nghiêng khuôn mặt. & Loại bỏ \\ \hline
LBPH & Sử dụng toán tử phân mẫu nhị phân cục bộ để bóc tách phân vùng cấu trúc ảnh. & Chống chịu tốt với thay đổi cường độ sáng / Độ chính xác giảm sâu khi khoảng cách chụp xa. & Loại bỏ \\ \hline
\textbf{ArcFace} & Hàm mất mát phạt góc Additive Angular Margin Loss tối ưu khoảng cách lớp~\cite{deng2019arcface}. & Độ chính xác cao vượt trội (99.77\% trên LFW)~\cite{deng2019arcface}, trích xuất vector 512 chiều ổn định. & \textbf{AI Server} \\ \hline
\textbf{FaceNet} & Hàm mất mát Triplet Loss tối ưu hóa khoảng cách giữa các cặp bộ ba mẫu ảnh~\cite{facenet2015}. & Hỗ trợ tốt định dạng tệp tin TensorFlow Lite~\cite{tensorflow_lite2021} thuần Python không yêu cầu compiler. & \textbf{App Client} \\ \hline
\end{tabular}
\end{table}

\textbf{Lý do lựa chọn:} Đề tài quyết định tích hợp thư viện chuyên sâu \textit{InsightFace}~\cite{insightface_github} dựa trên mô hình ArcFace~\cite{deng2019arcface} cho môi trường máy chủ AI Production do thuật toán này sở hữu khả năng phân tách biên đặc trưng mạnh mẽ nhất hiện nay. Song song đó, thư viện \textit{FaceNet-PyTorch}~\cite{facenet2015} được cấu hình làm giải pháp hỗ trợ trên môi trường hệ điều hành Windows và app di động nhằm loại bỏ hoàn toàn các rào cản lỗi biên dịch mã nguồn C++ phức tạp của các công cụ nền tảng khác, bảo đảm tính ổn định thực tế cho luồng đối soát có cấu hình ngưỡng đạt 0.65.

\clearpage

\section{Căn Cước Công Dân gắn chip và giao thức NFC}
\label{section:3.2}

\subsection{Nền tảng lý thuyết bảo mật tài liệu quốc gia}

Thẻ Căn Cước Công Dân gắn chip điện tử do Bộ Công an Việt Nam phát hành tuân thủ nghiêm ngặt các quy định kỹ thuật quốc tế của tiêu chuẩn ICAO Doc 9303~\cite{icao9303} dành cho các loại tài liệu du lịch có thể đọc được bằng máy. Chip mã hóa tích hợp trên thẻ tổ chức lưu trữ các tệp phân vùng dữ liệu theo kiến trúc Cấu trúc dữ liệu logic LDS. 

Trong đó, hệ thống khai thác tệp dữ liệu DG2 lưu trữ thực thể ảnh chân dung gốc chất lượng cao do cơ quan quản lý thu chụp trực tiếp để làm ảnh chuẩn gốc tham chiếu, loại trừ việc tổ chức phải tự thu thập sinh trắc học thủ công. Kênh giao tiếp không dây trường gần NFC bắt buộc phải vượt qua bước bắt tay kiểm soát truy cập cơ bản BAC sử dụng cặp khóa phiên sinh từ vùng máy đọc MRZ tại mặt sau của thẻ nhằm triệt tiêu các nguy cơ tấn công đọc trộm dữ liệu từ xa.

\subsection{Biện luận giải pháp kỹ thuật lựa chọn}

Giao thức này được lựa chọn nhằm hiện thực hóa yêu cầu bóc tách thông tin nhân thân chính thống tại thực địa vạch ra tại phân đoạn mục 2.1. Đề tài thực hiện so sánh đối chiếu giữa hai phương thức kết nối đọc thẻ thông dụng qua bảng ma trận kỹ thuật.

\begin{table}[ht!]
\centering
\small
\renewcommand{\arraystretch}{1.3}
\caption{Bảng đối chiếu phương thức trích xuất dữ liệu thẻ nhân sự}
\label{tab:nfc_comparison}
\begin{tabular}{|>{\raggedright\arraybackslash}p{2.3cm}|>{\raggedright\arraybackslash}p{4.5cm}|>{\raggedright\arraybackslash}p{5.5cm}|>{\centering\arraybackslash}p{1.7cm}|}
\hline
\textbf{Phương thức} & \textbf{Hạ tầng phần cứng} & \textbf{Đặc trưng kỹ thuật / Bảo mật} & \textbf{Lựa chọn} \\ \hline
Mạch quẹt thẻ RFID/Mifare & Sử dụng đầu đọc thẻ từ tần số 13.56 MHz kết nối nối tiếp qua cổng Serial. & Chi phí rẻ, tốc độ nhanh / Dữ liệu mã ID đơn giản, không có chữ ký số bảo chứng, dễ làm giả. & Loại bỏ \\ \hline
\textbf{Đọc chip NFC\newline ICAO 9303} & Tận dụng bộ cảm biến sóng ngắn NFC tích hợp sẵn trên điện thoại di động trạm Android. & Trích xuất ảnh gốc DG2 mã hóa sâu, dữ liệu được bảo chứng bằng chữ ký số SOD quốc gia. & \textbf{Lựa chọn\newline duy nhất} \\ \hline
\end{tabular}
\end{table}

\textbf{Lý do lựa chọn:} Đề tài quyết định ứng dụng thư viện nội bộ \texttt{cccd\_vietnam} xây dựng trên nền tảng giao thức \texttt{dmrtd} viết bằng ngôn ngữ Dart của framework Flutter. Giải pháp này cho phép ứng dụng di động trạm đầu cuối thực hiện chuỗi mã lệnh bắt tay an toàn BAC một cách tự động, giải mã luồng dữ liệu nhị phân bóc tách từ phân vùng tệp tin DG2 trực tiếp thành định dạng hình ảnh chân dung để truyền gửi kèm vào Request Body phục vụ cho tiến trình đối soát chéo lớp 2 của máy chủ trung tâm.

\clearpage

\section{Giao thức MQTT và mô hình Publish/Subscribe}
\label{section:3.3}

\subsection{Nền tảng lý thuyết truyền thông mạng IoT}

MQTT (\textit{Message Queuing Telemetry Transport})~\cite{mqtt5} là giao thức truyền thông điệp mạng theo kiến trúc gọn nhẹ, hoạt động dựa trên mô hình Publish/Subscribe, cực kỳ tối ưu cho các hệ thống phần cứng điều khiển ngoại vi có tài nguyên bộ nhớ hạn chế và vận hành trong môi trường hạ tầng mạng không dây chập chờn, không ổn định. Toàn bộ luồng truyền dữ liệu giữa thực thể phát và thực thể nhận đều được định tuyến tường minh thông qua máy chủ môi giới trung gian MQTT Broker~\cite{hivemq} dựa trên cấu trúc phân cấp cây của các Topics. Giao thức hỗ trợ ba cấp độ quản lý chất lượng dịch vụ bảo đảm dữ liệu truyền bao gồm QoS 0, QoS 1 và QoS 2.

\subsection{Biện luận giải pháp kỹ thuật lựa chọn}

Giao thức MQTT over SSL/TLS được lựa chọn nhằm xử lý yêu cầu phát lệnh điều khiển rơ-le đóng ngắt khóa điện từ từ máy chủ xuống vi điều khiển phần cứng ngoại vi dưới 500ms vạch ra tại phân đoạn mục 2.1. Đề tài phân tích so sánh kỹ thuật với giao thức truyền thông thay thế.

\begin{table}[ht!]
\centering
\small
\renewcommand{\arraystretch}{1.3}
\caption{Bảng so sánh giao thức điều khiển thiết bị phần cứng ngoại vi vật lý}
\label{tab:mqtt_comparison}
\begin{tabular}{|>{\raggedright\arraybackslash}p{2.3cm}|>{\raggedright\arraybackslash}p{4.5cm}|>{\raggedright\arraybackslash}p{5.5cm}|>{\centering\arraybackslash}p{1.7cm}|}
\hline
\textbf{Giao thức} & \textbf{Cơ chế hoạt động} & \textbf{Độ trễ hệ thống / Băng thông mạng} & \textbf{Lựa chọn} \\ \hline
HTTP Polling & Vi điều khiển ESP32 chạy vòng lặp gửi request liên tục lên server sau mỗi 1 giây. & Độ trễ lớn ($\ge 1000$ms), gây lãng phí tài nguyên máy chủ và chiếm dụng băng thông rộng. & Loại bỏ \\ \hline
\textbf{MQTT MQTTS} & Thiết lập một kết nối TCP duy nhất kéo dài, Broker chủ động đẩy gói lệnh tức thời. & Độ trễ cực thấp (120ms - 280ms), gói tin tiêu đề tinh giản tối đa chỉ vài byte dữ liệu. & \textbf{Lựa chọn\newline duy nhất} \\ \hline
\end{tabular}
\end{table}

\textbf{Lý do lựa chọn:} Sử dụng hạ tầng điện toán đám mây \textit{HiveMQ Cloud} làm Broker trung tâm kết nối qua giao thức bảo mật mã hóa SSL. Giao thức MQTT triệt tiêu hoàn toàn gánh nặng xử lý của HTTP truyền thống, giúp bo mạch vi điều khiển ESP32~\cite{espressif2023} thu nhận lệnh mở cửa \texttt{"action":"OPEN\_DOOR"} gần như tức thời ngay khi máy chủ Backend phê duyệt kết quả điểm danh thành công. Nhằm tối ưu hóa ô nhiễm tiếng ồn thực địa và bám sát triết lý thiết kế cảnh báo trực quan, hạ tầng MQTT được cấu hình tinh giản tuyệt đối nhiệm vụ, chỉ tập trung đóng ngắt rơ-le vật lý, giải phóng hoàn toàn các mô-đun phần cứng âm thanh/loa kéo tại hiện trường.

\clearpage

\section{Kiến trúc hướng dịch vụ và kiến trúc REST API}
\label{section:3.4}

\subsection{Nền tảng lý thuyết thiết kế phần mềm phân tán}

Kiến trúc hướng dịch vụ (Service-Oriented Architecture - SOA) là phương pháp luận thiết kế hệ thống phần mềm phân tán bằng cách phân rã ứng dụng thành các dịch vụ thành phần độc lập, sở hữu ranh giới nghiệp vụ rõ ràng và giao tiếp chéo thông qua môi trường mạng internet. Chuyển dịch sang mô hình REST~\cite{fielding2000}, các tài nguyên dữ liệu của hệ thống được định vị tường minh thông qua các địa chỉ URI và thao tác xử lý thông qua hệ thống phương thức HTTP tiêu chuẩn theo nguyên lý phi trạng thái, đóng gói payload định dạng cấu trúc JSON gọn nhẹ.

\subsection{Biểu luận giải pháp kỹ thuật lựa chọn}

Định hướng kiến trúc này giải quyết trực tiếp yêu cầu thiết lập mô hình hệ thống phân tán, đa nền tảng và có năng lực cô lập lỗi cho từng module xử lý vạch ra tại phân hệ kiến trúc mục 2.2. Đề tài thực hiện so sánh đối chiếu với mô hình kiến trúc khối cứng truyền thống.

\begin{table}[ht!]
\centering
\small
\renewcommand{\arraystretch}{1.3}
\caption{Bảng đối chiếu so sánh mô hình kiến trúc phát triển hệ thống}
\label{tab:architecture_comparison}
\begin{tabular}{|>{\raggedright\arraybackslash}p{2.3cm}|>{\raggedright\arraybackslash}p{4.5cm}|>{\raggedright\arraybackslash}p{5.5cm}|>{\centering\arraybackslash}p{1.7cm}|}
\hline
\textbf{Kiến trúc} & \textbf{Đặc trưng kết khối} & \textbf{Khả năng bảo trì / Khả năng mở rộng} & \textbf{Lựa chọn} \\ \hline
Monolithic & Toàn bộ mã nguồn Java, UI, và các mô hình AI gom chung vào một file thực thi duy nhất. & Triển khai nhanh ban đầu / Xung đột thư viện hệ thống, crash một module gây sập toàn bộ ứng dụng. & Loại bỏ \\ \hline
\textbf{REST API\newline (Phân tán)} & Phân tách độc lập hoàn toàn giữa Backend, Web Frontend và AI Service riêng biệt. & Dễ dàng nâng cấp, thay thế mã nguồn Face Engine AI mà không ảnh hưởng tới lõi nghiệp vụ. & \textbf{Lựa chọn\newline duy nhất} \\ \hline
\end{tabular}
\end{table}

\textbf{Lựa chọn giải pháp:} Việc tách biệt phân hệ xử lý sinh trắc học sang một dịch vụ AI Python Flask độc lập giúp hệ thống bảo vệ lõi nghiệp vụ Java Spring Boot an toàn, tối ưu hóa được năng lực xử lý bất đồng bộ đa luồng cấu hình tài nguyên của máy chủ, bám sát đặc tả kỹ thuật tiêu chuẩn tài liệu OpenAPI 3.0.

\clearpage

\section{JSON Web Token và phân quyền dựa trên vai trò}
\label{section:3.5}

\subsection{Nền tảng lý thuyết an ninh bảo mật hạ tầng mạng}

JSON Web Token (JWT) là một tiêu chuẩn xác thực mở quốc tế (chuẩn RFC 7519)~\cite{rfc7519} quy định quy trình đóng gói và truyền tải thông tin an toàn giữa các bên dưới dạng một đối tượng JSON nhỏ gọn được ký số bảo mật. Cấu trúc mã token JWT bao gồm ba phân đoạn độc lập phân tách bởi dấu chấm: Header, Payload và Signature. Cơ chế này kết hợp chặt chẽ với mô hình kiểm soát truy cập dựa trên vai trò - RBAC, phân định đặc quyền hệ thống theo danh mục phân quyền tăng dần rõ ràng.

\subsection{Biện luận giải pháp kỹ thuật lựa chọn}

Giải pháp mã hóa phi trạng thái này được lựa chọn nhằm thỏa mãn yêu cầu bảo mật thông tin an toàn API, phòng chống tấn công giả mạo quyền truy cập vạch ra tại phân đoạn mục 2.1 and 2.4. Đề tài lập bảng đánh giá kỹ thuật với giải pháp Session truyền thống.

\begin{table}[ht!]
\centering
\small
\renewcommand{\arraystretch}{1.3}
\caption{Bảng so sánh giải pháp quản lý phiên xác thực an ninh API}
\label{tab:security_comparison}
\begin{tabular}{|>{\raggedright\arraybackslash}p{2.3cm}|>{\raggedright\arraybackslash}p{4.5cm}|>{\raggedright\arraybackslash}p{5.5cm}|>{\centering\arraybackslash}p{1.7cm}|}
\hline
\textbf{Giải pháp} & \textbf{Cơ chế lưu trữ phiên} & \textbf{Năng lực chịu tải / Tính linh hoạt mạng} & \textbf{Lựa chọn} \\ \hline
Session Cookie & Máy chủ lưu trữ trạng thái phiên trong bộ nhớ RAM tập trung (HttpSession). & Gây quá tải bộ nhớ máy chủ khi số lượng kết nối tăng cao, khó mở rộng mô hình phân tán. & Loại bỏ \\ \hline
\textbf{JWT Token\newline + RBAC} & Phi trạng thái (Stateless), Token tự chứa thông tin xác thực ký số an toàn. & Máy chủ không tốn tài nguyên lưu trữ, dễ dàng kiểm tra quyền qua thuật toán Signature. & \textbf{Lựa chọn\newline duy nhất} \\ \hline
\end{tabular}
\end{table}

\textbf{Lý do lựa chọn:} Ứng dụng mã JWT giúp máy chủ Spring Boot lập tức giải mã xác chuẩn danh tính người dùng web và định danh thiết bị trạm đầu cuối ngoại vi tức thời trong từng request gửi lên thông qua chuỗi cấu hình Header \texttt{Authorization: Bearer} mà không cần thực hiện các câu lệnh truy vấn tốn kém tài nguyên vào cơ sở dữ liệu MySQL cốt lõi. Từng endpoint xử lý nghiệp vụ được bảo vệ nghiêm ngặt ở mức mã nguồn bằng các bộ lọc an ninh kết hợp chặt chẽ với Spring Security.

\clearpage

\section{WebSocket và cập nhật dữ liệu thời gian thực}
\label{section:3.6}

\subsection{Nền tảng lý thuyết đường ống song công}

WebSocket là giao thức truyền thông song công toàn phần hoạt động thông qua một kết nối TCP duy nhất kéo dài bền vững, được chuẩn hóa quốc tế bởi tổ chức IETF theo văn bản cấu hình RFC 6455~\cite{rfc6455}. Giao thức loại bỏ hoàn toàn cơ chế Request-Response thụ động của HTTP truyền thống bằng cách thiết lập một kênh kết nối lý tưởng kéo dài, cho phép máy chủ trung tâm chủ động đẩy trực tiếp các gói dữ liệu trạng thái xuống các client đang lắng nghe ngay khi sự kiện xuất hiện tại thực địa mà không cần đợi yêu cầu từ phía trình duyệt giao diện người dùng.

\subsection{Biện luận giải pháp kỹ thuật lựa chọn}

Giao thức WebSocket được lựa chọn nhằm hiện thực hóa yêu cầu tự động cập nhật lịch sử điểm danh trực tuyến lên màn hình giám sát hiển thị Web Dashboard tức thời mà không cần thực hiện thao tác tải lại trang vạch ra tại phân đoạn mục 2.1. Đề tài so sánh với giải pháp Polling.

\begin{table}[ht!]
\centering
\small
\renewcommand{\arraystretch}{1.3}
\caption{Bảng đối chiếu giải pháp đồng bộ dữ liệu giao diện người dùng Web Dashboard}
\label{tab:websocket_comparison}
\begin{tabular}{|>{\raggedright\arraybackslash}p{2.3cm}|>{\raggedright\arraybackslash}p{4.5cm}|>{\raggedright\arraybackslash}p{5.5cm}|>{\centering\arraybackslash}p{1.7cm}|}
\hline
\textbf{Giải pháp} & \textbf{Cơ chế đồng bộ dữ liệu} & \textbf{Độ trễ dữ liệu / Tài nguyên máy chủ} & \textbf{Lựa chọn} \\ \hline
AJAX Polling & Trình duyệt React cài đặt hàm Interval tự động gửi request gọi API sau mỗi 2 giây. & Độ trễ dữ liệu lớn, tạo ra hàng loạt request thừa gây nghẽn mạch cục bộ hệ thống máy chủ. & Loại bỏ \\ \hline
\textbf{WebSocket\newline STOMP} & Thiết lập kênh kết nối kéo dài độc lập, Server chủ động phát broadcast sự kiện. & Dữ liệu hiển thị trực tuyến tức thời (dưới 1 giây), tối ưu hóa băng thông mạng truyền tải. & \textbf{Lựa chọn\newline duy nhất} \\ \hline
\end{tabular}
\end{table}

\textbf{Lý do lựa chọn:} Lõi Backend Spring Boot kết hợp chặt chẽ với giao thức nhắn tin hướng văn bản đơn giản STOMP thông qua phân hệ cấu hình bộ định tuyến. Phía ứng dụng client React đăng ký lắng nghe ổn định trên chủ đề mạng cục bộ \texttt{/topic/attendance} vận hành chuyên biệt tại cổng mạng 3002, giúp toàn bộ dữ liệu chỉ số điểm danh hiển thị trực quan không đồng bộ mượt mà dưới 1 giây. Đường ống này gánh vác toàn bộ hạ tầng Visual Alert, tự động phân loại và mã hóa các biến cố vi phạm chuỗi trạng thái thành các Badge thông báo màu cam như \texttt{MISSING\_IN\_AUTO\_OPEN} hoặc \texttt{ANTI\_PASSBACK\_IN} để hỗ trợ OPERATOR giám sát từ xa một cách trực quan, loại bỏ các thông báo còi hú làm phiền không gian làm việc.

\clearpage

\section{Các framework và công nghệ tích hợp sử dụng}
\label{section:3.7}

Để triển khai hiện thực hóa thành công mô hình kiến trúc phân tán đa nền tảng đáp ứng đầy đủ các yêu cầu thông số phi chức năng về hiệu năng tốc độ xử lý luồng đồng thời của hệ thống vạch ra tại mục 2.4, đồ án lựa chọn tích hợp các công nghệ lõi tiên tiến kết nối đa tầng:

\begin{figure}[ht!]
    \centering
    \includegraphics[width=0.90\linewidth]{Hinhve/KienTrucHongDichVu.png}
    \caption{Sơ đồ cấu trúc kiến trúc tích hợp và điều phối luồng dữ liệu hệ thống}
    \label{fig:architecture_total_system}
\end{figure}

\begin{itemize}
    \item \textbf{Spring Boot 3.2:}~\cite{springboot3} Framework Java mã nguồn mở phiên bản v3.2.3 đóng vai trò quản lý lớp nghiệp vụ máy chủ trung tâm, cung cấp cơ chế tự động cấu hình tích hợp sẵn máy chủ nhúng Embedded Tomcat Server hỗ trợ kiến trúc Multi-Port phân tách tài nguyên - cổng 8080 cho lõi nghiệp vụ và 8081 cho luồng truyền lệnh phần cứng ngoại vi MQTT, bao bọc hạ tầng liên kết luồng bất đồng bộ để phục vụ hiển thị Web Dashboard React tại cổng mạng 3002, kết hợp thư viện Spring Data JPA quản lý tầng ánh xạ đối tượng và Spring Security thực thi bộ lọc an ninh bảo mật bám sát tiêu chuẩn OWASP Top 10.
    \item \textbf{React 18:}~\cite{react18} Thư viện JavaScript mã nguồn mở phiên bản v18.2.0 vận hành dựa trên cơ chế quản lý DOM ảo tối ưu kết hợp mô hình phân rã cấu phần giao diện, phối hợp cùng HTTP Client Axios và thư viện Chart.js hiển thị trực quan các biểu đồ thống kê nhật ký lịch sử điểm danh tại cổng 3002.
    \item \textbf{Flutter:}~\cite{flutterdart} Framework phát triển ứng dụng di động đa nền tảng của Google sử dụng ngôn ngữ lập trình Dart, sở hữu năng lực biên dịch mã nguồn trực tiếp xuống ngôn ngữ máy giúp tối ưu hóa hiệu năng phần cứng trạm đầu cuối di động Android, gọi mượt mọc các API hệ thống đọc chip kết nối trường gần và xử lý hiển thị banner đồ họa nhắc nhở chuỗi trực quan trên màn hình.
    \item \textbf{MySQL 8.0:}~\cite{mysql2021} Hệ quản trị cơ sở dữ liệu quan hệ mã nguồn mở mạnh mẽ, tích hợp bộ lưu trữ InnoDB hỗ trợ toàn vẹn giao dịch ACID nghiêm ngặt, cấu hình kiểu dữ liệu JSON chuyên dụng cho phép lưu trữ mảng chuỗi dữ liệu vector nhúng đặc trưng 512 chiều khuôn mặt cá nhân, quản lý hồ chứa kết nối qua kết nối pool của thư viện HikariCP tối ưu hóa tài nguyên.
\end{itemize}

Thảo luận tổng kết chương, Chương 3 đã phân tích chi tiết hệ thống nền tảng lý thuyết toán học của không gian vector đặc trưng khuôn mặt Cosine similarity, cấu trúc tệp dữ liệu bảo mật chip CCCD theo chuẩn ICAO Doc 9303, giao thức mạng IoT MQTT, kiến trúc hướng dịch vụ REST API, cơ chế bảo mật xác thực Token JWT và kênh truyền song công toàn phần WebSocket. Các bài toán đối chiếu so sánh biện luận khoa học chỉ rõ lý do lựa chọn giải pháp kỹ thuật thay thế đặt ra trong phần Tổng quan đều đã được làm rõ dứt khoát. 

Từ các cơ sở lý thuyết khoa học vững chắc và tập hợp các framework công nghệ đã được định hình tường minh này, tiến trình bóc tách thiết kế chi tiết cấu trúc mã nguồn, cài đặt API thực tế cho từng module phần mềm Backend Spring Boot, dịch vụ AI Flask, giao diện React Web Dashboard và cấu hình firmware định tuyến điều khiển thiết bị chấp hành ngoại vi sẽ được trình bày cụ thể trong chương tiếp theo — Chương 4.

\end{document}